\clearpage
\Article{358.\ 359}
\textsc{Theorem.} \emph{Let $P$ be a prime function and $x^\nu$ the lowest power of $x$ which, divided by $P$, gives remainder $1$. Further, let $P\equiv(P,\rho^n)$. Then $n$ is congruent to some power of the number $p$ modulo $\nu$.}

\emph{Proof.} We showed above that, if $P$
\[
=x^m+Ax^{m-1}+Bx^{m-2}+\text{etc.}
\]
then
\[
\begin{aligned}
z^m+Az^{m-1}+Bz^{m-2}+\text{etc.}\qquad\qquad\qquad\qquad&\\
{}-(z-x)(z-x^p)\dotstwo(z-x^{p^{m-1}})&
\end{aligned}
\]
is divisible by $P$. In the same way it would follow that
\[
\begin{aligned}
z^m+Az^{m-1}+Bz^{m-3}+\text{etc.}\qquad\qquad\qquad\qquad&\\
{}-(z-x^n)(z-x^{np})\dotstwo(z-x^{np^{m-1}})&
\end{aligned}
\]
is divisible by $P$. Since these factors are coprime to one another, the individual factors must necessarily be congruent to the corresponding individual factors modulo $P,p$. Thus $z-x^n$ must be $\equiv z-x^{p^\kappa}$, that is, $p^\kappa\equiv n(\mathrm{mod.}\,\nu)$. Q. E. D.\begingroup\interfootnotelinepenalty=10000\relax\renewcommand{\thefootnote}{*)}\footnote{If $(P,\rho^a)\equiv(P,\rho^b)(\mathrm{mod.}\,p)$, then $a\equiv p^\kappa b(\mathrm{mod.}\,\nu)$.

\emph{Proof.} Let $z^m+Az^{m-1}+Bz^{m-2}+\dotstwo=\Pi$. Then $(\Pi,\rho^a)\equiv(\Pi,\rho^b)(\mathrm{mod.}\,P)$; but
\[
\begin{aligned}
(\Pi,\rho^a)&\equiv(z-x^a)(z-x^{ap})(z-x^{app})\dotstwo(z-x^{ap^{m-1}}),\\
(\Pi,\rho^b)&\equiv(z-x^b)(z-x^{bp})(z-x^{bpp})\dotstwo(z-x^{bpm-1})
\end{aligned}
\]
which establishes the proposition.

The product of $\Pi,(\Pi,\rho^2),(\Pi,\rho^3)$, and so on, $(\Pi,\rho^\nu)$ is $\equiv(z^\nu-1)^m(\mathrm{mod.}\,P)$, for
\[
\begin{aligned}
(z-x)(z-x^2)(z-x^3)\ldots(z-x^\nu)&\\
{}\equiv(z-x^p)(z-x^{2p})(z-x^{3p})\dotstwo(z-x^{\nu p})&\equiv\text{etc.}\equiv z^\nu-1
\end{aligned}
\]
In the sequence $P,(P,\rho^2),(P,\rho^3)$, and so on, $\dotsfour(P,\rho^\nu)$, all prime divisors of the function $x^\nu-1$ occur, each $m$ times. It follows that the product of them all is $\equiv(x^\nu-1)^m$.}\endgroup

\par\medskip\Needspace{5\baselineskip}
\begin{center}\emph{Finding the prime divisors of the function $x^\nu-1$ for a prime modulus.}\end{center}
\Article{360}
If $\nu$ is divisible by the modulus $p$ or by one of its powers, put $\nu=p^k\lambda$. Then
\[
x^\nu-1\equiv(x^\lambda-1)^{p^k}(\mathrm{mod.}\,p)
\]
This makes it clear that only the case where $\nu$ is not divisible by $p$ needs to be considered.

\SourceLeaf{240}
If $p^m\equiv1(\mathrm{mod.}\,\nu)$, with $m$ as small as possible, then $x^{p^m-1}-1$ is clearly divisible by $x^\nu-1$. Thus $x^\nu-1$ cannot have any divisors other than those of $x^{p^m-1}-1$. But this expression has prime divisors of degree $m$ and others whose degrees divide the number $m$. Thus $x^\nu-1$ also has divisors of these kinds. We shall show how many it has of each kind by an example, from which the general law can easily be deduced.

Let $\nu=63$ and $p=13$. Then $m=6$. Thus $x^{63}-1$ has prime factors of degrees six, three, two and one modulo $13$. It is clear that the product of the factors of degree one is the common divisor of greatest degree of the functions $x^{63}-1$ and $x^{12}-1$, that is, $x^3-1$; hence there are three prime factors of degree one. The product of all prime factors of degrees two and one is the common divisor of the functions $x^{63}-1$ and $x^{168}-1$, that is, $x^{21}-1$; hence there are $\frac{21-3}2$, or $9$, factors of degree two. The product of the prime factors of degrees three and one is the common divisor of the functions $x^{63}-1$ and $x^{2196}-1$, that is, $x^9-1$; hence there are $\frac{9-3}3$, that is, $2$, divisors of degree three. Finally, the remaining ones have degree six, and their number is therefore $=\frac{63-6-18-3}6$, that is, $6$.

\Needspace{12\baselineskip}
Careful consideration of this matter easily gives the following general rule:

\emph{Let $\delta$ be a divisor of $m$, and let all divisors of the number $\delta$ less than $\delta$ be $\delta',\delta'',\delta'''$, and so on. Let the greatest common divisors of $\nu$ with $p^\delta-1$, $p^{\delta'}-1$, $p^{\delta''}-1$, and so on, be $\mu,\mu',\mu''$, and so on, respectively. Put $\frac\mu{\mu'};\frac\mu{\mu''},\frac\mu{\mu'''}$, and so on, $=\lambda',\lambda'',\lambda'''$, and so on. Then $x^\nu-1$ has $\frac1\delta$ times as many prime divisors of degree $\delta$ as there are numbers below $\mu$ divisible by none of the numbers $\lambda',\lambda'',\lambda'''$, and so on.}

\Needspace{10\baselineskip}
\Article{361}
\textsc{Theorem.} \emph{If a function $X$ of the indeterminate $x$ is divisible by another function $\xi$, and $X$, when $x$ is replaced by $x^k$, becomes $X'$, then $X'$ is divisible by $(\xi,\rho^{\frac1k})$.}

\emph{Proof.} Let $X\equiv\xi\upsilon$, and let $\xi,\upsilon$ become $\xi',\upsilon'$ when $x$ is replaced by $x^k$. It is clear that $X'\equiv\xi'\upsilon'$. But $\xi'$ is divisible by $(\xi,\rho^{\frac1k})$. Hence so is $X'$. \mbox{Q. E. D.}
